% =========================================================================
% Glossary (unnumbered, front matter, roman pages). Two pages maximum.
% =========================================================================

\section*{Glossary}
\addcontentsline{toc}{section}{Glossary}

\begingroup\footnotesize\singlespacing   % scoped: an unscoped size switch here shrank the whole body once
Terms specific to this thesis, or given a specific sense here; ordinary
machine-learning vocabulary is not repeated. Chapter and section numbers
point to a term's fullest treatment, not necessarily its first use.

\begin{description}[leftmargin=1.8em,labelsep=0.4em,itemsep=0.5pt,parsep=0pt,topsep=2pt]

\item[ancestor cone (cone)] The over-approximate upstream subgraph of $U$
that could have shaped a failure; not a counterfactual, and never a claim
of causation on its own.

\item[archaeology] The method of re-reading this project's own earlier
campaign archives to check whether a candidate's ``discovery'' already
existed there. One of the three substitute endpoints for H1.

\item[band (same-config band)] The range of round-to-round score movement
produced by re-measurement noise alone, under an unchanged configuration; a
move inside it carries no information.

\item[capability vs efficacy] The distinction, drawn out in the clinic
(\S\ref{sec:clinic}), between whether the harness \emph{can} take an
action (capability) and whether taking it \emph{changes an outcome}
(efficacy). The \emph{capability wall} is the rival account that the
action space, not the readout, cannot reach the relevant lesion.

\item[candidate] A proposed edit to the harness's configuration, tools,
prompt or processors, authored by the Evolver, before it faces the gates.

\item[the clinic] This thesis's name, by extension of its drug/dose/trial
vocabulary, for the capability-vs-efficacy testing regime of
\S\ref{sec:clinic}.

\item[composition graph ($G$)] A typed, static re-expression of the
harness's configuration: nodes are hook points, processors, slots,
bundles, skills and tools.

\item[Critic] The meta-role that judges each candidate, may interrogate
the Evolver, and audits the candidate set against the scoreboard before
the gates run.

\item[deaf (selection is deaf)] Said of an instrument whose noise floor
exceeds the magnitude of the effect it is meant to detect, so a null
reading cannot distinguish ``no effect'' from ``undetectable effect.''

\item[Digester] The meta-role that reads every trajectory, pass and fail
alike, and writes a per-task dossier.

\item[dossier] The Digester's per-task write-up of a trajectory: what
happened, and the reusable strategy or fragility it names.

\item[Evolver] The meta-role that authors zero or more candidate edits,
each as a manifest plus an applied configuration.

\item[\texttt{\textbackslash F\{n\}} / \cls{A}/\cls{B}/\cls{C}] The anchor
notation: every number in this thesis carries an \F{}~tag into the
numbered row of Appendix~\ref{app:ledger} that supports it; the row is
classed \cls{A} (recomputable by a named script), \cls{B} (recorded, not
regenerable) or \cls{C} (judgment).

\item[flip / flip rate] A task whose pass/fail outcome differs between two
rounds run under an identical configuration; the flip rate is the fraction
of the bed that does this (\S\ref{sec:what-graph-did}).

\item[four layers of mortality] This thesis's taxonomy of the four reasons
a correctly-prescribed edit fails to produce a lasting gain
(Chapter~\ref{ch:results}).

\item[genotype / deployment / phenotype] Three nested identity hashes for
a candidate: its structural configuration; that configuration plus its
runtime overlay; and the two plus the execution edges actually observed.

\item[harness] Everything around the model that turns it into an agent:
tools, prompts, processors and the control flow that calls them. ``An LLM
agent is a model plus a harness.''

\item[level 2] Three unrelated senses collide in this thesis: GAIA
difficulty tier~2 (a benchmark property, always written in that form);
rung~2 of the readout ladder (``fired, and on what''); and the official
cheatsheet's Level-1/Level-2 \emph{capability evidence} that every
candidate must carry (the artifact loads; the mechanism round-trips),
which Appendix~\ref{app:deviations} calls capability evidence throughout.

\item[lift] A localization metric: the hit rate on a shipped candidate's
named target tasks, divided by the base rate of those tasks flipping
anyway.

\item[no-pixel bed] The 100-task subset of GAIA's text-only split used for
every number in this thesis, excluding the three tasks whose content the
recorder cannot attribute.

\item[Planner] The meta-role that synthesizes every dossier in a round
into one landscape document for the Evolver.

\item[positive control] The one intervention in this study known, on
mechanism grounds, to have worked; used to calibrate what the score axis
can detect.

\item[processor] A pluggable unit of harness behaviour, composed into
bundles and hung on hook points; the object the composition graph types.

\item[ratchet] The paper's specified --- and, in the released code,
unimplemented --- rule that a task solved in one round may never be given
back in a later one.

\item[readout / readout ladder (L0--L3)] The evidence available about a
candidate edit, in four rungs: \textbf{L0} deployed, \textbf{L1} executed,
\textbf{L2} fired (and on what), \textbf{L3} changed outcomes.
``Readout'' is this evidence generally.

\item[recorder] The write-only component that logs the composition graph
and $U$ during a run, without feeding back into the loop's own decisions.

\item[resolution] In this thesis, the smallest score change
distinguishable from same-configuration noise (a detection-limit sense) ---
not GAIA's own integer scoring resolution, and not the VIM instrumentation
sense.

\item[same-configuration window] A run, or pair of runs, executed under
one unchanged configuration, used to measure noise rather than an edit's
effect.

\item[seam] A point in the official loop's six stages at which GHX
attaches from outside, without altering what runs there by default.

\item[ship / \texttt{no\_op}] The two terminal decisions a round can
reach: land one or more candidates (\emph{ship}), or change nothing
(\emph{no\_op}).

\item[sixth gate] GHX's added acceptance gate: a candidate's claimed
target node must appear in the recorded execution graph of the failure it
cites, or the candidate is refused.

\item[survivor (the survivor)] The one positive channel in this study ---
episodic notes prepended to the prompt --- that outlasted every governance
and dormancy failure mode the rest of the loop suffered.

\item[treadmill (re-prescription treadmill)] The pattern by which the
same edit family is independently re-derived across campaigns without
inheriting from an earlier instance that already worked.

\item[typed candidate surface] GHX's seven typed graph-edit operations,
replacing free-text candidate authoring with edits a build-time validator
can check.

\item[unfolded execution graph ($U$)] The per-invocation DAG of what a run
actually executed, with edges observed rather than reconstructed or
inferred.

\item[whitelist / whitelist campaign] The evidence-admissibility ruling
(2026-08-26) restricting every claim in this thesis to specific campaigns;
a \emph{whitelist campaign} is one admitted under it.

\end{description}
\endgroup
